\chapter{Instrumental Variables}
\label{cap:iv}
% ── index ───────────────────────────────────────────────────────
\index{instrumental variables|textbf}
\index{IV|see{instrumental variables}}
\index{iv@\texttt{iv()}|textbf}
\index{endogeneity|textbf}
\index{two-stage least squares}
\index{2SLS|see{two-stage least squares}}
\index{relevance condition|textbf}
\index{exclusion restriction|textbf}
\index{first stage}

% ─────────────────────────────────────────────────────────────────────────────
\section{The endogeneity problem}

Consider the structural model:
\[
  y = X_1\beta_1 + X_2\beta_2 + \varepsilon
\]

where $X_1$ are exogenous regressors ($n \times k_1$) and $X_2$ are potentially
endogenous regressors ($n \times k_2$). OLS is biased and inconsistent when:
\[
  E[\varepsilon \mid X_2] \neq 0
  \quad\Longleftrightarrow\quad
  \mathrm{plim}\,\hat{\beta}_{\mathrm{OLS}} \neq \beta
\]

The most common causes are:

\begin{itemize}
  \item \textbf{Omitted variable} correlated with $X_2$: if $y = X\beta + W\gamma + u$
        and $W$ is omitted, then $\varepsilon = W\gamma + u$ and $\mathrm{Cov}(X_2, \varepsilon) \neq 0$.
  \item \textbf{Simultaneity}: $y$ and $X_2$ are determined jointly.
  \item \textbf{Classical measurement error}: $X_2 = X_2^* + \nu$, with $X_2^*$
        the true value; the attenuation factor is $\mathrm{plim}\,\hat{\beta} = \beta \cdot \sigma_{X^*}^2/\sigma_X^2 < \beta$.
\end{itemize}

% ─────────────────────────────────────────────────────────────────────────────
\section{Identification conditions}

Let $Z$ be the $n \times L$ matrix of instruments (including $X_1$ as
included instruments). A valid instrument satisfies two conditions:

\begin{enumerate}
  \item \textbf{Relevance:}
        $\mathrm{rank}\bigl(E[z_i x_i']\bigr) = k$, where $k = k_1 + k_2$.
        In practice: $X_2$ is correlated with $Z$ in the first stage.

  \item \textbf{Exogeneity (exclusion restriction):}
        $E[z_i \varepsilon_i] = 0$.
        The excluded instruments affect $y$ \emph{only} through $X_2$.
\end{enumerate}

The model is identified when $L \geq k$ (order condition). For
$L = k$ (exact identification) the IV estimator is unique; for $L > k$
(overidentification) 2SLS is efficient in the class of IV estimators.

% ─────────────────────────────────────────────────────────────────────────────
\section{The 2SLS estimator}

Define the projection matrix onto the column space of $Z$:
\[
  P_Z = Z(Z'Z)^{-1}Z'
\]

The Two-Stage Least Squares (2SLS) estimator is:
\[
  \boxed{\hat{\beta}_{\mathrm{2SLS}} = (X'P_Z X)^{-1} X'P_Z y}
\]

\subsection*{Two-stage interpretation}

\textbf{First stage} --- projects each column of $X_2$ onto $Z$:
\[
  X_2 = Z\hat{\Pi} + \hat{V}, \qquad \hat{\Pi} = (Z'Z)^{-1}Z'X_2
\]

\textbf{Second stage} --- regresses $y$ on $\hat{X} = P_Z X = [X_1,\, \hat{X}_2]$:
\[
  \hat{\beta}_{\mathrm{2SLS}} = (\hat{X}'X)^{-1}\hat{X}'y
\]

\begin{aviso}
  Never run the two stages manually as sequential OLS. The standard errors
  from the second stage will be wrong because $\hat{X}_2$ is estimated.
  Always use \lstinline{iv()}, which corrects the variance automatically.
\end{aviso}

% ─────────────────────────────────────────────────────────────────────────────
\section{Asymptotic variance}

Under exogeneity and relevance, $\hat{\beta}_{\mathrm{2SLS}}$ is consistent and
asymptotically normal:
\[
  \sqrt{n}\,(\hat{\beta}_{\mathrm{2SLS}} - \beta)
  \;\xrightarrow{d}\;
  \mathcal{N}\!\bigl(0,\; \sigma^2 \,Q_{XZ}^{-1} Q_{ZZ} Q_{XZ}^{-1\prime}\bigr)
\]

where $Q_{XZ} = E[x_i z_i']$ and $Q_{ZZ} = E[z_i z_i']$.

The variance estimator is:
\[
  \widehat{\mathrm{Var}}[\hat{\beta}_{\mathrm{2SLS}}]
  = s^2_{\mathrm{IV}}\,(X'P_Z X)^{-1},
  \qquad
  s^2_{\mathrm{IV}} = \frac{\hat{e}'_{\mathrm{2SLS}}\,\hat{e}_{\mathrm{2SLS}}}{n - k}
\]

The same sandwich corrections from OLS (HC1, HC3, clustered) apply
by replacing $X$ with $\hat{X} = P_Z X$.

% ─────────────────────────────────────────────────────────────────────────────
\section{Weak instruments}

An instrument is weak when its correlation with $X_2$ is low. The bias of
2SLS relative to OLS is approximately:
\[
  \frac{\mathrm{Bias}_{\mathrm{2SLS}}}{\mathrm{Bias}_{\mathrm{OLS}}}
  \approx \frac{1}{F + 1}
\]

where $F$ is the $F$-statistic from the first stage. For $F < 10$ (the rule of
thumb from Staiger \& Stock, 1997) the instrument is considered weak.

The \textbf{Cragg-Donald statistic} generalizes to multiple endogenous
regressors:
\[
  CD = \frac{n - L}{L - k_1} \cdot \lambda_{\min}
\]

where $\lambda_{\min}$ is the smallest eigenvalue of the concentration matrix.
Stock \& Yogo (2005) tabulate critical values for maximum relative bias
of 5\%, 10\%, 20\%, and 30\% relative to OLS.

% ─────────────────────────────────────────────────────────────────────────────
\section{DGP Verification}

The same method as Chapter~\ref{cap:ols} applies to IV: a DGP with known
endogeneity is defined, OLS and IV are estimated, and it is verified that IV
corrects the bias.

\subsection*{DGP structure}

Let $z$ and $v$ be independent random variables $\sim \mathcal{N}(0,1)$.
Define $\varepsilon$ as the \emph{residual} from projecting $v$ onto $[1,\;z]$:
\[
  \varepsilon = v - \delta_0 - \delta_1 z,
  \qquad
  \delta_1 = \frac{\mathrm{Cov}_n(z,v)}{\mathrm{Var}_n(z)},
  \quad
  \delta_0 = \bar v - \delta_1\bar z
\]

By construction $\sum \varepsilon_i = 0$ and $\sum z_i\varepsilon_i = 0$, i.e.,
$Z'\varepsilon = 0$ \emph{exactly} in sample. Then:

\[
  x = 0{,}8\,z + \varepsilon
  \qquad
  y = 1{,}5 + 2{,}0\,x + 0{,}7\,\varepsilon
\]

The structural error $0{,}7\varepsilon$ creates endogeneity ($\mathrm{Cov}(x,0{,}7\varepsilon) \neq 0$),
but since $Z'\varepsilon = 0$ exactly, the IV moment conditions are
satisfied without sampling noise. Result: IV recovers $(\beta_0, \beta_1)$
to machine precision --- the same as OLS did without noise in Chapter~\ref{cap:ols}.

The OLS bias is predictable:
\[
  \mathrm{plim}\;\hat\beta_{\mathrm{OLS}}
  = \beta_1 + \frac{\mathrm{Cov}(x,\varepsilon)}{\mathrm{Var}(x)}
  \approx 2{,}0 + \frac{0{,}7}{0{,}64 + 1}
  \approx 2{,}43
\]

\subsection*{Script}

\begin{lstlisting}[caption={DGP with endogeneity --- OLS vs.\ IV}]
seed(42)
input df
id
1
end
for i in 1..=10 { df = append(df, df) }
df |> filter(_n <= 1000)
generate df z   = rnormal()
generate df v   = rnormal()
let d1          = cov(df, z, v) / variance(df, z)
let d0          = mean(df, v) - d1 * mean(df, z)
generate df eps = v - d0 - d1*z
generate df x   = 0.8*z + eps
generate df y   = 1.5 + 2.0*x + 0.7*eps
let m_ols = ols(y ~ x, df)
let m_iv  = iv(y ~ x, ~ z, df)
esttab(m_ols, m_iv)
\end{lstlisting}

\noindent Actual interpreter output:

\begin{lstlisting}[language={}, basicstyle=\ttfamily\small,
  frame=none, numbers=none, backgroundcolor=\color{codbg},
  xleftmargin=0pt, xrightmargin=0pt, breaklines=false]
                            (1)              (2)
                            OLS          IV/2SLS
────────────────────────────────────────────────
x                     2.4303***        2.0000***
                       (0.0108)         (0.0280)
const                                  1.5000***
                                        (0.0129)
Constant              1.3273***
                       (0.0058)
────────────────────────────────────────────────
N                          1000             1000
R²                       0.9807
Adj. R²                  0.9807
────────────────────────────────────────────────
* p<0.10  ** p<0.05  *** p<0.01
\end{lstlisting}

IV recovers $(\hat\beta_0, \hat\beta_1) = (1{,}5000,\; 2{,}0000)$ to machine
precision. OLS overestimates $\beta_1$ by $+0{,}43$ (2.4303 vs.\ 2.0000),
consistent with the predicted asymptotic bias of $+0{,}43$.
$R^2$ is not reported for IV since it has no standard interpretation under endogeneity.

% ─────────────────────────────────────────────────────────────────────────────
\section{Syntax}

\begin{lstlisting}
iv(structural_formula, instrument_formula, df)
iv(structural_formula, instrument_formula, df, cov=type)
\end{lstlisting}

The \textit{structural formula} specifies the equation of interest. The
\textit{instrument formula} lists the excluded instruments after \lstinline{~}:

\begin{lstlisting}[caption={Basic IV --- one instrument}]
// structural equation: y ~ x  (x is endogenous)
// excluded instrument: z
iv(y ~ x, ~ z, df)
\end{lstlisting}

For multiple endogenous regressors and instruments:

\begin{lstlisting}
iv(y ~ x1 + x2, ~ z1 + z2, df)
\end{lstlisting}

\subsection{Transformations and interactions in formulas}
\label{subsec:iv-formula-advanced}

The same transformations and interactions available in OLS
(Section~\ref{subsec:formula-transformations}--\ref{subsec:formula-I})
work equally in the structural and instrument formulas of IV:

\begin{lstlisting}[caption={IV with transformations and interactions}]
// transformed endogenous regressor
iv(Y ~ log(X2), ~ Z, df)

// interaction of transformations in the structural equation
iv(Y ~ log(K) + log(L) + log(K):log(L), ~ Z1 + Z2, df)

// transformed instrument
iv(Y ~ X2, ~ log(Z), df)
\end{lstlisting}

\begin{nota}
  Included exogenous regressors ($X_1$) do not need to be listed in the
  instrument formula --- they are added automatically as their own
  instruments.
\end{nota}

\section{Standard error options}

The same \lstinline{cov=} options available in OLS apply to IV:

\begin{lstlisting}
iv(y ~ x, ~ z, df)                           // classical
iv(y ~ x, ~ z, df, cov=robust)              // HC1
iv(y ~ x, ~ z, df, cov=HC3)                 // HC3
iv(y ~ x, ~ z, df, cov=cluster, cluster=id) // clustered
\end{lstlisting}

\section{Capturing the result}

\begin{lstlisting}
let m_iv = iv(y ~ x, ~ z, df)
display m_iv
\end{lstlisting}

\section{Weak instrument diagnostics}

\begin{lstlisting}[caption={Weak instrument test}]
weak_iv(y ~ x, ~ z, df)
\end{lstlisting}

The function reports:
\begin{itemize}
  \item First-stage $F$-statistic (per endogenous variable)
  \item Cragg-Donald statistic $CD$
  \item Stock \& Yogo (2005) critical values for 5\%--30\% relative bias
  \item Reference: $F > 10$ (Staiger \& Stock, 1997)
\end{itemize}

\section{IV post-estimation tests}

\subsection{Sargan / Hansen J --- overidentification}

When there are more instruments than endogenous regressors
(overidentification), test the exogeneity of the instruments:

\begin{lstlisting}
estat_overid(y ~ x1 + x_endog, ~ z1 + z2, df)
\end{lstlisting}

$H_0$: instruments are exogenous (valid). The Sargan statistic is
$n \cdot R^2$ from the regression of IV residuals on the instruments,
distributed as $\chi^2(L - K)$, where $L$ is the number of instruments
and $K$ the number of regressors. Rejecting $H_0$ suggests the
instruments may be invalid. Not applicable when the model is exactly
identified ($L = K$).

\subsection{Durbin-Wu-Hausman --- endogeneity}

Tests whether the suspect variable is truly endogenous, comparing OLS
and IV:

\begin{lstlisting}
estat_endog(y ~ x1 + x_endog, ~ z1 + z2, df)
\end{lstlisting}

$H_0$: regressors are exogenous (OLS is consistent). Uses the augmented
regression approach: adds first-stage residuals to the structural
equation and tests their significance via $F$. If $H_0$ is rejected, IV
is needed; if not rejected, OLS is consistent and preferred (more
efficient).

\section{Fitted values}

\begin{lstlisting}
let m_iv = iv(y ~ x, ~ z, df)
// $\hat{y} = X\hat{\beta}_{\mathrm{2SLS}}$
let yhat = predict(m_iv, "xb")
\end{lstlisting}

\section{Comparison table}

The typical workflow quantifies the endogeneity bias by comparing OLS and IV:

\begin{lstlisting}[caption={Complete IV analysis workflow}]
load "data.csv" as df

let m_ols = ols(wage ~ educ, df)
let m_iv  = iv(wage ~ educ, ~ distance, df, cov=robust)

// first-stage diagnostics
weak_iv(wage ~ educ, ~ distance, df)

// comparison table
esttab(m_ols, m_iv)
\end{lstlisting}

If $\hat{\beta}^{\mathrm{IV}} > \hat{\beta}^{\mathrm{OLS}}$, the OLS bias
is negative --- which is consistent with attenuation from measurement error or with
an omitted variable that depresses both $X_2$ and $y$.
